\documentclass[11pt]{article}
\usepackage{fontspec}
\usepackage[french]{babel}
\usepackage{amsmath}
\usepackage{amssymb}
\usepackage{amsthm}
\usepackage{mathrsfs}
\usepackage{tikz-cd}
\usepackage{enumitem}

\newcommand{\printedpage}[1]{\marginpar{\footnotesize p.~#1}}
\newcommand{\msp}[1]{\marginpar{\footnotesize [#1]}}
% marker of a manuscript page standing inside a displayed array (where \marginpar is not possible)
\newcommand{\mspin}[1]{\texttt{[#1]}}
\newcommand{\HOT}{\mathrm{HOT}}
\newcommand{\Cat}{\mathrm{Cat}}
\newcommand{\Fl}{\mathrm{Fl}}
\newcommand{\Ob}{\mathrm{Ob}}
\newcommand{\Ens}{\mathrm{Ens}}
\newcommand{\Hom}{\underline{\mathrm{Hom}}}
\newcommand{\defeq}{\overset{\text{déf}}{=}}
\newcommand{\pl}{\texttt{[plutôt}}
\newcommand{\uHom}{\underline{\operatorname{Hom}}}

\begin{document}

\printedpage{24}
% head: A. Grothendieck : Dérivateurs, Ch. II, Édition des Universités de Montpellier II et Paris VII
% foot: Version 8 juillet 2001
\msp{page 37}

\section*{3\quad Retour sur la cofinalité}

Il est instructif de relier les questions de cofinalité (pour une flèche $u : I' \longrightarrow I$ dans Cat) au dérivateur $\HOT_{W_0}$, où $W_0 \subseteq \Fl\,\Cat$ est le plus grossier des localiseurs dans Cat satisfaisant l'axiome de connexité Loc 5, i.e.
\[
W_0 = \{ f \in \Fl\,\Cat \mid \pi_0(f) \text{ bijectif}\}. \tag*{(3.1)}
\]
On trouve alors, pour $I$ dans Cat, une équivalence canonique, 2-fonctorielle en $I$,
\[
\HOT_{W_0}(I) \xrightarrow{\;\sim\;} I^{\wedge} \defeq \uHom(I,(\Ens)), \quad \text{\pl}\ \uHom(I^o,\Ens)\,\text{\texttt{]}} \tag*{(3.2)}
\]
en d'autres termes, on a une équivalence
\[
\HOT_{W_0} \simeq \mathbf{D}_{(\Ens)}. \tag*{(3.3)}
\]
Quand $\HOT_{W_0}$ est défini par localisation à partir de la catégorie $\Cat(I) \defeq \uHom(I^o,\Cat)$, la flèche (3.2) est celle déduite, par localisation, de la flèche
\[
\begin{array}{rcl}
\pi_0^I : \uHom(I^o,\Cat) & \longrightarrow & \uHom(I^o,(\Ens)) \\[1ex]
X & \longmapsto & \pi_0 \circ X,
\end{array} \tag*{(3.4)}
\]
où
\[
\pi_0 : \Cat \longrightarrow (\Ens) \tag*{(3.5)}
\]
est le foncteur ‘composantes connexes’. Si on décrit $\HOT_{W_0}(I)$ comme localisée de $\Cat/I$, (3.2) est déduit de la flèche
\[
\left\{
\begin{array}{rcl}
\Cat/I & \longrightarrow & I^{\wedge} = \uHom(I^o,(\Ens)) \\[1ex]
X & \longmapsto & \underline{\pi_0}(X/I) = (i \longmapsto \pi_0(i\backslash X)).
\end{array}
\right. \tag*{(3.6)}
\]

\msp{page 38}
Ainsi, les foncteurs familiers, pour une flèche
\[
u : I' \longrightarrow I
\]
dans Cat
\[
\left\{
\; I'^{\wedge} \quad
\begin{array}{c}
\xrightarrow{\;u_!\;} \\[1.5ex]
\xleftarrow{\;u^*\;} \\[1.5ex]
\xrightarrow{\;u_*\;}
\end{array}
\quad I^{\wedge},
\right. \tag*{(3.7)}
\]
\printedpage{25}
% head: A. Grothendieck : Dérivateurs, Ch. II, Édition des Universités de Montpellier II et Paris VII
% foot: Version 8 juillet 2001
s'interprètent comme les foncteurs de même nom pour le dérivateur $\HOT_{W_0}$, ou plus directement, pour le dérivateur (qu'on pourrait faussement considérer comme canularesque) $\mathbf{D}_{(\Ens)}$.

Cela donne un sens à la cohorte des notions cohomologiques associées à un dérivateur, ou plus particulièrement, à un localiseur fondamental $W$ sur Cat (donnant naissance au dérivateur $\HOT_W$). On les désignera par le nom de $W_0$-fourbi (au lieu de $\HOT_{W_0}$-fourbi, ou $\mathbf{D}_{(\Ens)}$-fourbi), notamment pour une flèche $u : I' \longrightarrow I$ de Cat telle la notion
\[
\left\{
\begin{array}{lll}
\text{$W_0$-équivalence} & \Bigl(\iff & \text{$\pi_0(u)$ bijectif.}\Bigr) \\
\text{$W_0$-asphérique} & \Bigl(\iff & \text{les $I'/i$ pour $i \in \Ob I$ sont} \\
 & & \text{$0$-connexes.}\Bigr) \\
\text{$W_0$-coasphérique} & \Bigl(\iff & \text{les $i\backslash I'$ pour $i \in \Ob I$ sont} \\
 & & \text{$0$-connexes.}\Bigr) \\
\text{$W_0$-propre (\footnotemark[46])} & \Bigl(\iff & \text{les inclusions $I'_i \longrightarrow I'/i$ ($i \in \Ob I$)} \\
 & & \text{sont $W_0$-coasphériques, i.e. les fibres} \\
 & & \text{des colocalisées $i'\backslash I' \longrightarrow i\backslash I$} \\
 & & \text{($i' = u(i)$) \pl\ ($i = u(i')$)]} \\
 & & \text{sont $0$-connexes.}\Bigr) \\
\text{\mspin{page 39}} & & \\
\text{$W_0$-lisse (\footnotemark[47])} & \Bigl(\iff & \text{les inclusions $I'_i \longrightarrow i\backslash I'$} \\
 & & \text{sont $W_0$-asphériques, i.e. les fibres} \\
 & & \text{des flèches localisées $I'/i' \longrightarrow I/i$} \\
 & & \text{($i = u(i')$) sont $0$-connexes.}\Bigr) \\
\text{$W_0$-parfait} & \Bigl(\iff & \text{$W_0$-propre et $W_0$-lisse.}\Bigr) \\
\text{$W_0$-équivalence universelle (\footnotemark[48])} & \Bigl(\iff & \text{$\pi_0$ bijectif et le reste par tout} \\
 & & \text{changement de base. Ce sont aussi} \\
 & & \text{les $W_0$-fibrations (ou $W_0$-fibrations} \\
 & & \text{faibles à fibres $W_0$-asphériques,} \\
 & & \text{i.e. $0$-connexes.}\Bigr) \\
\text{$W_0$-fibration} & \Bigl(\iff & \text{$u^* : \Cat/I \longrightarrow \Cat/I'$ transforme} \\
 & & \text{$W_0$-équivalences en $W_0$-équivalences.}\Bigr) \\
\text{$W_0$-fibration faible (\footnotemark[49])} & \Bigl(\iff & \text{le changement de base} \\
 & & \text{$u^* : \Cat/I \longrightarrow \Cat/I'$ transforme} \\
 & & \text{homotopismes $J_1 \longrightarrow J_2$ en} \\
 & & \text{$W_0$-équivalences - et il suffit de} \\
 & & \text{l'exiger pour l'inclusion d'un objet} \\
 & & \text{final (ou au choix initial) et pour} \\
 & & \text{l'inclusion d'un point dans $\boldsymbol{\Delta}^1$.}\Bigr)
\end{array}
\right. \tag*{(3.8)}
\]
\footnotetext[46]{Exemple : les Cat-cofibrations.}
\footnotetext[47]{Exemple : les Cat-fibrations.}
\footnotetext[48]{Exemple : les flèches $W_0$-lisses ou $W_0$-propres, à fibres 0-connexes.}
\footnotetext[49]{Exemple : les flèches $W_0$-lisses telles que $\pi_0(I'/I) = (i \longmapsto \pi_0(i\backslash I')) = u_!(e_{I'^{\wedge}})$ soit localement constant sur $I$, i.e. $\underline{\pi_0}(I'/I)$ soit un \textit{revêtement} étale de $I$, ou si $u^{\mathrm{opp}}$ satisfait cette condition.}

\printedpage{26}
% head: A. Grothendieck : Dérivateurs, Ch. II, Édition des Universités de Montpellier II et Paris VII
% foot: Version 8 juillet 2001
\texttt{[Nous avons une]} implication \textit{stricte} (et comment)
\[
\text{$W_0$-équivalence universelle} \Longrightarrow \text{$W_0$-fibration}
\]
\texttt{[et une]} implication \textit{stricte} (et comment\,!) (\footnote[50]{Mais si les fibres sont 0-connexes, les deux notions de $W_0$-fibrations sont équivalentes et signifient qu'on a une $W_0$-équivalence universelle \dots})
\[
\text{$W_0$-fibration} \Longrightarrow \text{$W_0$-fibration faible.}
\]

Signalons, comme exemple d'utilisation de ces notions, que pour un diagramme cartésien dans Cat
\[
\begin{tikzcd}
I' \arrow[d, "u"'] & J' \arrow[l, "p'"'] \arrow[d, "v"] \\
I & J, \arrow[l, "p"']
\end{tikzcd} \tag*{(3.9)}
\]
la flèche de changement de base pour $u_*$, $v_*$ pour les préfaisceaux est \texttt{[un]} isomorphisme si $u$ est $W_0$-propre, ou $p$ $W_0$-lisse, et la flèche de changement de base pour $u_!$, $v_!$ est \texttt{[un]} isomorphisme si $u$ est $W_0$-lisse ou $p$ $W_0$ propre (\footnote[51]{Sauf erreur, ca donne \textit{deux caractérisations} des morphismes $W_0$-propres, ou des morphismes $W_0$-lisses.}). Je signale à ce propos que ces énoncés, qui a priori se formulent pour les préfaisceaux à valeurs dans (Ens), restent valables % [sic] ca -> ça (in footnote 51), missing cedilla
\msp{page 40}
pour des préfaisceaux à valeurs dans une catégorie $\mathcal{M}$ quelconque (sous réserve d'existence des operations d'images directes cohomologiques ou homologiques envisagées). Cela se voit p.ex. en notant que % [sic] operations -> opérations, missing accent
\[
\text{$W_0$-équivalence} \;\Longrightarrow\; \text{$\mathbf{D}_{\mathcal{M}}$-équivalence} \tag*{(3.10)}
\]
pour toute $\mathcal{U}$-catégorie $\mathcal{M}$, d'où résulte que toute (\footnote[52]{Pas tout à fait exacte pour \textit{toute} ‘$W_0$-notion’, il faut excepter les deux notions de $W_0$-fibration forte, et faible.}) ‘$W_0$-notion’ implique la $\mathbf{D}_{\mathcal{M}}$-notion correspondante, et permet donc d'appliquer à $\mathbf{D}_{\mathcal{M}}$ les résultats le concernant. (Il y a bien sûr un argument direct, moins savant, consistant à se ramener au cas où $\mathcal{M}$ est petite, et à utiliser alors le plongement pleinement fidèle $\mathcal{M} \hookrightarrow \mathcal{M}^{\wedge}$ resp. $\mathcal{M} \hookrightarrow ((\mathcal{M}^o)^{\wedge})^o$.)

Je rappelle pour mémoire une (troisième\,?) caractérisation des flèches $u : I' \longrightarrow I$ $W_0$-propres (resp. $W_0$-lisses) : ce sont celles pour lesquelles le foncteur changement de base $u^* : \Cat/I \longrightarrow \Cat/I'$ transforme flèches $W_0$-coasphériques (resp. $W_0$-asphériques) en itou. Cela gagne en intérêt, du fait que la notion de foncteur $W_0$-coasphérique (resp. $W_0$-asphérique) s'interprète comme celle de foncteur cofinal (resp. co-cofinal), cf. plus bas.

\printedpage{27}
% head: A. Grothendieck : Dérivateurs, Ch. II, Édition des Universités de Montpellier II et Paris VII
% foot: Version 8 juillet 2001
\msp{page 41}
La notion de $W_0$-fibration apparaît comme extrêmement forte - beaucoup plus forte à certains égards que celle (disons) de $W_\infty$-fibration ou de $W_\omega$-fibration. Elle équivaut à ceci (comme on voit aussitôt) : pour toute composante connexe $I'_\alpha$ de $I'$, désignant par $I_\alpha$ la composante connexe correspondante de $I$, et pour tout $J_\alpha$ 0-connexe et au dessus de $I_\alpha$, le produit fibré $I'_\alpha \times_{I_\alpha} J_\alpha$ est encore 0-connexe. (Ça ressemble aux histoires de totale 0-connexité d'un topos\,!) Mais cela signifie aussi que $I'_\alpha \longrightarrow I_\alpha$ est une $W_0$-équivalence universelle. Cela implique que ses fibres sont 0-connexes, et la réciproque est vraie si $u : I' \longrightarrow I$ est $W_0$-propre ou $W_0$-lisse. Ainsi, la condition essentielle pour une $W_0$-fibration est celle que les fibres des flèches induites $I'_\alpha \longrightarrow I_{u(\alpha)}$ \pl\ $I'_\alpha \longrightarrow I_\alpha$\texttt{]} soient 0-connexes. En fait, on a ceci

\noindent\textbf{Proposition 3.1.} \textit{Soit $u : I' \longrightarrow I$ une flèche dans} Cat.
\begin{itemize}
\item[\textit{a)}] \textit{Pour que ce soit une $W_0$-équivalence universelle, il faut et il suffit que pour toute $\boldsymbol{\Delta}^1 \longrightarrow I$, i.e. toute flèche $\alpha$ de $I$, la catégorie $I' \times_I \boldsymbol{\Delta}^1 \defeq I'|\alpha$ soit $0$-connexe, ou encore que les fibres de $I$ soient $0$-connexes, et que pour toute flèche} non identique \textit{$\alpha$ de $I$, $I'|\alpha$ soit $0$-connexe.}
\item[\textit{b)}] \textit{Pour que $u : I' \longrightarrow I$ soit une $W_0$-fibration, il faut et il suffit que pour toute composante connexe $I'_{\alpha'}$ de}
\msp{page 42}
\textit{$I'$, posant $\alpha = u(\alpha')$, la flèche induite $I'_{\alpha'} \longrightarrow I_\alpha$ soit une $W_0$-équivalence universelle (cf. critère a) pour cette notion).}
\item[\textit{c)}] \textit{Pour que ce soit une $W_0$-fibration faible, il} faut \textit{que $\forall\, i \in I$, l'inclusion $I'_i \hookrightarrow i\backslash I'$ soit une $W_0$-équivalence, i.e. $\pi_0(I'_i) \xrightarrow{\;\sim\;} \pi_0(i\backslash I')$ soit bijectif,} et \textit{que $i \longmapsto \pi_0(i\backslash I')$ soit un préfaisceau localement constant sur $I$ (i.e. $\underline{\pi_0}(I'/I)$ soit un} revêtement \textit{étale de $I$). Cette condition est aussi suffisante, si $u$ est $W_0$-lisse. De plus (pour mémoire) $u$ est une $W_0$-fibration faible si et seulement si $u^{\mathrm{opp}}$ l'est.}
\end{itemize}

La plus grande différence entre la notion de $W_0$-fibration, et celle de $W_0$-fibration faible, me paraît que les morphismes $W_0$-parfaits (i.e. $W_0$-lisses et $W_0$-propres), et a fortiori les $W$-parfaits pour un localiseur fondamental $W \subseteq W_0$, sont forcément des $W$-fibrations faibles, mais (\textit{même si $W = W_\omega$,} le plus fin de tous les localiseurs fondamentaux) ne sont pas des $W$-fibrations. L'exemple le plus frappant à cet égard est sans doute celui des revêtements étales, qui sont $W_\omega$-parfaits (donc
\msp{page 43}
des $W_\omega$-fibrations faibles et a fortiori des $W_0$-fibrations faibles), mais pas des $W_0$-fibrations\,; car dans le cas où $I'$, $I$ sont 0-connexes, il faudrait que les fibres soient 0-connexes (cf. b) ci-dessus), i.e. que $u$ soit un isomorphisme\,!

Une propriété sympathique des morphismes $W_0$-lisses (resp. $W_0$-propres), c'est que ce sont des applications \textit{ouvertes} (resp. fermées) pour les espaces topologiques sous-jacents, ce qui signifie en effet que les fibres locales, i.e. des flèches localisées en haut $I'/i' \longrightarrow I/i$ (resp. colocales, i.e. les flèches colocalisées en bas $i'\backslash I' \longrightarrow i\backslash I$), sont surjectives, i.e. à fibres non-vides (ce qui est impliqué par la condition que celles-ci soient 0-connexes).

\printedpage{28}
% head: A. Grothendieck : Dérivateurs, Ch. II, Édition des Universités de Montpellier II et Paris VII
% foot: Version 8 juillet 2001
Dans tout ça je n'ai encore guère parlé des flèches $W_0$-asphériques resp. co-asphériques, qui pourtant sont ici une motivation principale pour introduire $W_0$\,!

\noindent\textbf{Proposition 3.2.} \textit{Conditions équivalentes sur la flèche $u : I' \longrightarrow I$ dans} Cat :
\begin{itemize}
\item[a)] \textit{Pour toute $\mathcal{U}$-catégorie $\mathcal{M}$ et tout foncteur $F : I \longrightarrow \mathcal{M}$, la flèche canonique (sous réserve que les petites limites inductives soient représentables dans $\mathcal{M}$),}
\[
\varinjlim_{I'} u^*(F) \longrightarrow \varinjlim_{I} F \tag*{(3.2.1)}
\]
\msp{page 44}
\textit{est} \texttt{[un]} \textit{isomorphisme.}
\item[a$'$)] \textit{Pour toute $\mathcal{U}$-catégorie $\mathcal{M}$ où les petites limites projectives sont représentables et tout $F \in \mathcal{M}(I) = \uHom(I^o,\mathcal{M})$, i.e. tout préfaisceau sur $I$ à valeurs dans $\mathcal{M}$, la flèche canonique}
\[
\varprojlim_{I} F \longrightarrow \varprojlim_{I'} u^*(F) \tag*{(3.2.2)}
\]
\textit{est} \texttt{[un]} \textit{isomorphisme.}
\item[a$''$)] (\footnote[53]{C'est une condition sur les morphismes de topos, savoir que si $F$ \texttt{[est]} dans $\mathcal{X}^{\wedge}$, on a
\[
H^0(\mathcal{X},F) \xrightarrow{\;\sim\;} H^0(\mathcal{X}',u^*F).
\]
Il y a \texttt{[une]} notion similaire pour les $H^i$ ($i \leq n$), à dégager. Notion de morphismes $W_0$-coasphériques de topos\,! C'est dans la nature d'une condition de 0-connexité relative, duale de la notion $W_0$-asphérique, qui signifie simplement que pour $U$ 0-connexe dans $U^{\wedge}$ \pl\ $\mathcal{X}^{\wedge}$\texttt{]}, $f^*(U)$ \pl\ $u^*(U)$\texttt{]} est 0-connexe. De même, on a sans doute des notions (‘duales’ des précédentes) de $W_n$-coasphéricité (bien comprise en principe si $\mathcal{X}$ est localement $n$-asphérique) (La $W_n$-coasphéricité, étant celle évoquée plus haut.)}) \textit{Comme a$'$, avec $\mathcal{M} = (\Ens)$, i.e. $F \in \Ob I^{\wedge}$.}~(\footnote[54]{\textbf{NB} On peut introduire
\begin{itemize}
\item[a$'''$)] Comme a), mais avec $\mathcal{M} = (\Ens)$.
\item[a$^{IV}$)] Comme a), avec le foncteur canonique $I' \longrightarrow I'^{\wedge}$.
\end{itemize}}).
\item[b)] \textit{La flèche $u$ est $W_0$-coasphérique, i.e. pour tout $i \in I$, $i\backslash I'$ est $0$-connexe.}
\end{itemize}
L' équivalence de a, a$'$, a$''$ est $\pm$ tautologique en termes des définitions (cf. SGA 4, I, 8\dots). La forme a$''$ s'écrit, posant $\mathbf{D} = \HOT_{W_0}$, que $\forall\, F \in \mathbf{D}(I)$, on veut que
\[
\underbrace{H^{\bullet}_{\mathbf{D}}(I,F)}_{\defeq p_{I,*}(F)} \longrightarrow \underbrace{H^{\bullet}_{\mathbf{D}}(I',u^*F)}_{\defeq p_{I',*}(F)} \qquad \text{\pl}\quad \underbrace{H^{\bullet}_{\mathbf{D}}(I',u^*F)}_{\defeq p_{I',*}(u^*F)}\ \text{\texttt{]}} \tag*{(3.2.3)}
\]
($p_I : I \longrightarrow e$ et $p_{I'} : I' \longrightarrow e$ les foncteurs canoniques) soit \texttt{[un]} isomorphisme. On sait que cela équivaut à $u$ $\mathbf{D}$-coasphérique. La démonstration est un jeu amusant entre $u_!$, $u^*$, $u_*$ à la fois. Il faut exprimer que le Hom d'un objet $M$ de $\mathbf{D}(e)$ ($=$ (Ens) en l'occurence) dans la flèche (3.2.3) est bijective. Ici, il % [sic] occurence -> occurrence, French spelling

\end{document}
